\chapter{Método Proposto}\label{chp:proposta}

\section{Considerações iniciais}

Este capítulo descreve o método desenvolvido para a combinação dinâmica de previsões de séries temporais utilizando um agente baseado em grandes modelos de linguagem. O objetivo deste trabalho foi investigar e desenvolver um método de combinação de modelos para previsão de valores em séries temporais (\textit{forecasting}) que integrasse abordagens lineares e não-lineares a partir de um conjunto diverso de modelos, incluindo métodos estatísticos e de aprendizagem de máquina, e que explorasse distintas representações dos dados durante a etapa de treinamento, obtidas a partir das transformações \textit{wavelet} e de Fourier, de modo a aumentar a diversidade do conjunto de preditores combinados. Para a etapa de combinação, foi desenvolvido o CREST, um método que delega a um agente baseado em modelo de linguagem a decisão de qual estratégia de combinação empregar para cada série, em substituição a uma regra fixa de ponderação aplicada indistintamente a todas as séries de um conjunto de dados, com o objetivo de superar os resultados obtidos pelos modelos individuais e por outras estratégias de combinação, como as baseadas em ponderação

O restante deste capítulo está organizado da seguinte forma. A Seção~\ref{sec:formulacao} formaliza o problema de combinação de predições resolvido pelo método. A Seção~\ref{sec:crest} apresenta a visão geral do CREST e o fluxo completo de execução. A Seção~\ref{sec:preparacao} descreve o conhecimento construído em nível de conjunto de dados, que precede a atuação do agente. A Seção~\ref{sec:diagnostico} descreve como cada série é caracterizada antes da decisão de combinação. A Seção~\ref{sec:agente} detalha o agente combinador, seu catálogo de ações e seu ciclo de decisão. A Seção~\ref{sec:aplicacao} descreve como a estratégia decidida é aplicada e os mecanismos de rastreabilidade do método. Por fim, a Seção~\ref{sec:consideracoes-finais-proposta} encerra o capítulo com as considerações finais sobre a proposta.

%\section{Formulação do problema} \label{sec:formulacao}
%Seja $\mathcal{D} = \{X_1, X_2, \dots, X_N\}$ um conjunto de $N$ séries temporais univariadas e $\mathcal{M} = \{M_1, M_2, \dots, M_M\}$ um conjunto fixo de $M$ métodos preditivos. Para cada série $X_i$, esses métodos produzem previsões para um horizonte comum de $h$ observações. O conjunto pode incluir métodos estatísticos, de aprendizagem de máquina e ingênuos, bem como variantes aplicadas a representações transformadas dos dados. A composição do \textit{pool} de modelos constitui uma escolha de projeto, de modo que o CREST opera sobre as previsões fornecidas, sem exigir uma arquitetura específica dos preditores. Para permitir a combinação, as previsões devem corresponder aos mesmos instantes e estar expressas na mesma escala da variável-alvo.

%Para cada série $X_i$, cada preditor $M_j \in \mathcal{M}$ produz previsões em quatro janelas distintas: três janelas de validação, $w \in \{1, 2, 3\}$, para as quais os valores reais $y_i^{(w)}$ são conhecidos, e uma janela de teste subsequente às janelas de validação, para a qual os valores reais não são conhecidos no momento da decisão de combinação. Denota-se por $\hat{y}*{i,j}^{(w)}$ a previsão do preditor $M_j$ para a série $X_i$ na janela $w$, e por $\hat{y}*{i,j}^{(\mathrm{teste})}$ a previsão desse mesmo preditor para a janela de teste. Todas as janelas possuem horizonte $h$, de modo que $y_i^{(w)}, \hat{y}*{i,j}^{(w)}, \hat{y}*{i,j}^{(\mathrm{teste})} \in \mathbb{R}^{h}$. A janela de teste é posterior às janelas de validação, e as previsões de cada janela são produzidas utilizando apenas informações disponíveis antes de sua respectiva origem de previsão.


%O problema de combinação consiste em definir, para cada série $X_i$, uma função $f_i$ que combine as previsões de teste $\{\hat{y}_{i,1}^{(\text{teste})}, \dots, \hat{y}_{i,M}^{(\text{teste})}\}$ em uma única previsão final $\hat{y}_i^{(\text{combinada})}$, de modo a minimizar o erro em relação aos valores reais da janela de teste. A escolha de $f_i$ para cada série é fundamentada exclusivamente nas evidências disponíveis nas três janelas de validação, isto é, na relação entre $\hat{y}_{i,j}^{(w)}$ e $y_i^{(w)}$ para $w \in \{1,2,3\}$ e $M_j \in \mathcal{M}$, e nas características observáveis da própria série $X_i$, sem qualquer acesso aos valores reais da janela de teste.

%A formulação permite selecionar estratégias de combinação distintas para diferentes séries, sem impor uma única função a todo o conjunto $\mathcal{D}$. No protocolo considerado, essa seleção se apoia em três janelas de validação, que fornecem evidências sobre o desempenho dos preditores em um número limitado de períodos.

\section{Formulação do problema}
\label{sec:formulacao}

Seja $\mathcal{D} = \{X_i\}_{i=1}^{N}$ um conjunto de históricos de $N$ séries temporais univariadas, em que $X_i = (x_{i,t})_{t=1}^{T_i}$ contém as $T_i$ observações disponíveis da série $i$ antes da janela de teste que terá os seus valores preditos. O índice $i$ identifica a série, enquanto $t$ indica a posição temporal de uma observação nessa série. Seja ainda $\mathcal{M} = \{M_j\}_{j=1}^{M}$ um conjunto fixo de $M$ métodos preditivos (ou \textit{pool} de preditores), identificados pelo índice $j$. Todos os métodos produzem previsões para um horizonte comum de $h$ observações.

O conjunto de preditores $\mathcal{M}$ pode incluir métodos estatísticos, de aprendizagem de máquina e ingênuos, bem como variantes aplicadas a representações transformadas dos dados. A composição do \textit{pool} de preditores é uma escolha de projeto, de modo que o CREST opera sobre as previsões fornecidas, sem exigir uma arquitetura ou seleção específica dos preditores. Para permitir a combinação, as previsões devem corresponder aos mesmos instantes e escala da variável-alvo.

Para cada série $i$, consideram-se três janelas de validação $w \in \{1,2,3\}$, e uma janela de teste posterior a elas. Seja $\tau_i^{(w)}$ a origem de previsão da janela de validação $w$, definida como o índice da última observação anterior a essa janela, com
\[
    1 \leq \tau_i^{(1)} < \tau_i^{(2)} < \tau_i^{(3)}
    \quad \text{e} \quad
    \tau_i^{(w)} + h \leq T_i,
    \qquad w \in \{1,2,3\}.
\]

O vetor de valores reais da janela de validação $w$ é
\[
    \mathbf{y}_i^{(w)}
    =
    \bigl(y_{i,k}^{(w)}\bigr)_{k=1}^{h}
    \in \mathbb{R}^{h},
    \qquad
    y_{i,k}^{(w)}
    =
    x_{i,\tau_i^{(w)}+k},
\]
em que $k \in \{1,\dots,h\}$ identifica o passo no horizonte
de previsão. Denota-se por
\[
    \hat{\mathbf{y}}_{i,j}^{(w)}
    =
    \bigl(\hat{y}_{i,j,k}^{(w)}\bigr)_{k=1}^{h}
    \in \mathbb{R}^{h}
\]
o vetor de previsões produzido pelo método $M_j$ para essa janela, cuja componente $\hat{y}_{i,j,k}^{(w)}$ corresponde à previsão de $x_{i,\tau_i^{(w)}+k}$. Cada vetor de previsões é produzido utilizando apenas informações disponíveis até a respectiva origem $\tau_i^{(w)}$.

A janela de teste compreende os $h$ instantes imediatamente posteriores a $T_i$. Denotam-se por $x_{i,T_i+k}$, para $k \in \{1,\dots,h\}$, os valores futuros da série $i$,
que não integram o histórico disponível $X_i$ e são desconhecidos no momento da decisão de combinação. Esses valores compõem o vetor
\[
    \mathbf{y}_i^{(\mathrm{teste})}
    =
    \bigl(x_{i,T_i+k}\bigr)_{k=1}^{h}
    \in \mathbb{R}^{h}.
\]

O vetor de previsões do método $M_j$ para essa janela é
\[
    \hat{\mathbf{y}}_{i,j}^{(\mathrm{teste})}
    =
    \bigl(\hat{y}_{i,j,k}^{(\mathrm{teste})}\bigr)_{k=1}^{h}
    \in \mathbb{R}^{h},
\]
em que $\hat{y}_{i,j,k}^{(\mathrm{teste})}$ corresponde à previsão de $x_{i,T_i+k}$. Essas previsões utilizam apenas informações disponíveis até $T_i$.

O problema de combinação consiste em definir, para cada série $i$, uma função
\[
    f_i : (\mathbb{R}^{h})^{M} \longrightarrow \mathbb{R}^{h}
\]
que combine os $M$ vetores de previsões de teste em uma única previsão final:
\[
    \hat{\mathbf{y}}_i^{(\mathrm{comb})}
    =
    f_i\!\left(
        \hat{\mathbf{y}}_{i,1}^{(\mathrm{teste})},
        \ldots,
        \hat{\mathbf{y}}_{i,M}^{(\mathrm{teste})}
    \right).
\]

Busca-se reduzir o erro dessa previsão em relação a $\mathbf{y}_i^{(\mathrm{teste})}$. Como os valores reais de teste são desconhecidos no momento da decisão, a escolha de $f_i$ é fundamentada exclusivamente nas previsões e nos valores reais das três janelas de validação,
\[
    \left\{
        \left(
            \mathbf{y}_i^{(w)},
            \hat{\mathbf{y}}_{i,1}^{(w)},
            \ldots,
            \hat{\mathbf{y}}_{i,M}^{(w)}
        \right)
    \right\}_{w=1}^{3},
\]
e nas características observáveis do histórico $X_i$ de todas as séries do conjunto de dados $\mathcal{D}$, sem acesso aos valores reais da janela de teste.

A formulação permite selecionar estratégias de combinação distintas para diferentes séries, sem impor uma única função a todo o conjunto $\mathcal{D}$. No protocolo considerado, as três janelas de validação fornecem evidências sobre o desempenho dos preditores em um número limitado de períodos.

\section{Visão geral do CREST} \label{sec:crest}
O CREST (\textit{Cross-series Reasoning for Ensemble Strategy selecTion}) é o método desenvolvido para resolver o problema formulado na Seção~\ref{sec:formulacao}. Para cada série $i$, o CREST utiliza um agente baseado em modelo de linguagem para selecionar, entre estratégias de combinação previamente definidas, aquela que determina a função $f_i$ a ser aplicada
aos vetores de previsões de teste.

Essa seleção considera a possibilidade de diferentes séries se beneficiarem de estratégias distintas e a disponibilidade de apenas três janelas de validação por série, cada uma com $h$ observações. Nesse cenário, o agente não estima pesos por otimização direta sobre as janelas de validação. Sua decisão se apoia em evidências extraídas tanto do histórico $X_i$ e dos resultados de validação da série $i$ quanto de informações agregadas das demais séries de $\mathcal{D}$. 

%Em vez de definir uma única função de combinação $f$ aplicada a todas as séries de $\mathcal{D}$, como fazem métodos de \textit{meta-learning} que treinam uma regra fixa de ponderação a partir de um conjunto de treinamento, o CREST delega a definição de $f_i$, série a série, a um agente baseado em modelo de linguagem. Essa escolha decorre diretamente das duas restrições discutidas na Seção~\ref{sec:formulacao}: como a estratégia de combinação mais adequada varia entre séries e entre conjuntos de dados, e como cada série dispõe de apenas três pontos de observação para fundamentar essa escolha, o agente não estima pesos por otimização direta sobre essa amostra reduzida, mas sim raciocina sobre um conjunto de evidências mais amplo, construído tanto a partir da série individual quanto do conjunto de dados como um todo, para decidir qual entre um repertório de estratégias de combinação já existentes deve ser empregada.



A execução do CREST é dividida em duas macro-fases, ilustradas na Figura~\ref{fig:crest_arquitetura}. 

\begin{figure}[H]
    \centering
    \includegraphics[width=\textwidth]{images/proposta/crest_arquitetura_2.png}
    \caption[Visão geral da arquitetura proposta.]{Visão geral da arquitetura do CREST, dividida em uma fase de preparação por conjunto de dados e uma fase de decisão por série, na qual um agente itera sobre um catálogo de ações até aceitar uma estratégia de combinação, ou até que critérios de parada sejam atingidos.}
    \label{fig:crest_arquitetura}
\end{figure}

Na primeira macro-fase, que ocorre uma única vez por conjunto de dados $\mathcal{D}$ e não envolve nenhum modelo de linguagem, são construídas duas formas de conhecimento entre séries: $i)$ um meta-modelo que aprende a relação entre características das séries e o erro de cada preditor, e $ii)$ uma avaliação em dados de validação, de quais estratégias de combinação tendem a funcionar melhor naquele conjunto de dados especificamente. Essas duas formas de conhecimento, descritas na Seção~\ref{sec:preparacao}, ficam disponíveis para todas as séries de $\mathcal{D}$ antes que qualquer decisão individual seja tomada.

Na segunda macro-fase, executada de forma independente para cada série $i$, com histórico $X_i \in \mathcal{D}$, o agente decide qual estratégia de combinação empregar. Essa fase começa pela caracterização da série e pela avaliação do desempenho de cada preditor $M_j \in \mathcal{M}$ nas janelas de validação e pelo estabelecimento de um conjunto de estratégias de combinação simples já testadas, que servem tanto de piso de desempenho quanto de ponto de partida para o agente. Em seguida, o agente entra em um ciclo iterativo, descrito na Seção~\ref{sec:agente}, no qual propõe uma estratégia, verifica seu desempenho nas próprias janelas de validação, e decide se a aceita ou se propõe uma nova estratégia, até encerrar o ciclo por aceite, por estagnação do desempenho, ou por esgotamento do orçamento de iterações. A estratégia resultante desse processo, que corresponde sempre ao melhor desempenho observado ao longo de todo o ciclo, incluindo as estratégias simples testadas antes de o agente atuar, determina a função $f_i$, aplicada aos vetores de previsões de teste $\hat{\mathbf{y}}_{i,j}^{(\mathrm{teste})}$, para $j \in \{1,\dots,M\}$, produzindo a previsão final $\hat{\mathbf{y}}_i^{(\mathrm{comb})}$.

\subsection{Pseudocódigo}

A seguir, o Algoritmo~\ref{alg:crest} apresenta a visão geral do método
proposto, explicitando as entradas e saídas, a preparação em nível de conjunto
de dados, a caracterização de cada série, a avaliação das estratégias iniciais,
o ciclo de decisão do agente, os critérios de parada e a aplicação da
estratégia selecionada. O mecanismo de pontuação, comum às duas macro-fases, é
detalhado no Algoritmo~\ref{alg:pontuar}, e o diagnóstico interpretativo, no
Algoritmo~\ref{alg:diagnostico}. A notação empregada é a da formulação da
Seção~\ref{sec:formulacao}, complementada por: $\mathbf{z}_i$, o vetor das 26
características de $X_i$; $\mathcal{G}_{-i}$, o meta-modelo entre séries
ajustado sem a série $i$; $S_{\text{fix}}$, o conjunto de estratégias fixas
descrito na Seção~\ref{sec:preparacao}, formado pela média, mediana e DBA sobre
$\mathcal{M}$, pela média e média aparada sobre os $\kappa$ preditores mais
estáveis, para $\kappa \in \{5,7,9\}$, e, quando disponível, pela ponderação
definida pelo meta-modelo; $\mathcal{H}_i$, o histórico de tentativas da série
$i$, formado por pares $(s, \text{score}_i(s))$; $\mathcal{A}$, o catálogo de 24
ações descritas na Seção~\ref{sec:agente}, ao qual se soma o aceite, que não
integra o catálogo e encerra o ciclo de decisão; $\hat{\mathbf{Y}}_i^{(\mathrm{teste})}$,
a coleção das previsões de teste dos $M$ preditores; $T_i$, a avaliação do
conjunto de preditores na série $i$; $\mathcal{C}_i$, o diagnóstico
interpretativo; $K$, o orçamento de iterações; $P$, a paciência de estagnação;
e $L$, o limite de repetições para falhas de comunicação. Em cada iteração, o
agente responde exclusivamente com um \textit{Thought}, uma \textit{Action} e
uma \textit{Action Input}.

\begin{algorithm}[H]
\caption{CREST: visão geral do método proposto}
\label{alg:crest}
\Entrada{$\mathcal{D}=\{X_i\}_{i=1}^{N}$; $\mathcal{M}=\{M_j\}_{j=1}^{M}$;
         horizonte $h$; janelas de validação
         $\{(\mathbf{y}_i^{(w)}, \hat{\mathbf{y}}_{i,j}^{(w)})\}_{i,j,w}$ com
         $w \in \{1,2,3\}$; previsões de teste
         $\hat{\mathbf{Y}}_i^{(\mathrm{teste})}$; catálogo $\mathcal{A}$ com
         $|\mathcal{A}|=24$; orçamento $K$; paciência $P$; limite $L$}
\Saida{$\{\hat{\mathbf{y}}_i^{(\mathrm{comb})}\}_{i=1}^{N}$ e métricas de
       desempenho}

\For{$i \gets 1$ \KwTo $N$}{
  $\mathbf{z}_i \gets \Act{Perfil}(X_i)$
}
\For{$i \gets 1$ \KwTo $N$}{
  $\mathcal{G}_{-i} \gets \Act{AjustarMetaModelo}(\mathcal{D} \setminus \{X_i\})$
}
$S_{\text{fix}} \gets \{\text{média}, \text{mediana}, \text{DBA}\}$
\For{$\kappa \in \{5, 7, 9\}$}{
  $S_{\text{fix}} \gets S_{\text{fix}} \cup \{\text{média}_\kappa, \text{média aparada}_\kappa\}$
}
\If{$\mathcal{G}_{-i}$ está disponível}{
  $S_{\text{fix}} \gets S_{\text{fix}} \cup \{\text{ponderada}(\mathcal{G}_{-i})\}$
}
\For{$i \gets 1$ \KwTo $N$}{
  \ForEach{$s \in S_{\text{fix}}$}{
    $\text{score}_i(s) \gets \Act{Pontuar}(i, s)$
  }
}
\For{$i \gets 1$ \KwTo $N$}{
  \ForEach{$s \in S_{\text{fix}}$}{
    $\text{prior}_i(s) \gets \frac{1}{N-1}\sum_{i' \neq i} \text{score}_{i'}(s)$
  }
}
\For{$i \gets 1$ \KwTo $N$}{
  $T_i \gets \Act{AvaliarConjunto}(i)$
  $\mathcal{C}_i \gets \Act{Diagnosticar}(\mathbf{z}_i, T_i)$
  $\mathcal{H}_i \gets \{(s, \text{score}_i(s)) : s \in S_{\text{fix}}\}$
  $t \gets 1$; $q \gets 0$; $s_i^{\text{acc}} \gets \text{nulo}$
  \While{$t \le K$ \textbf{e} $q < P$ \textbf{e} $s_i^{\text{acc}} = \text{nulo}$}{
    $u \gets \Act{Contexto}(t, K, \mathbf{z}_i, T_i, \mathcal{C}_i, \text{prior}_i, \mathcal{H}_i)$
    $r \gets \Act{Agente}(u)$
    \While{($r$ vazia \textbf{ou} erro de conexão) \textbf{e} repetições $< L$}{
      $r \gets \Act{Agente}(u)$
    }
    $(\text{ação}, \text{argumentos}) \gets \Act{Validar}(r)$
    \eIf{formato inválido, ação fora do catálogo \textbf{ou} argumentos inválidos}{
      devolver erro estruturado ao agente
    }{
      \eIf{$\text{ação} = \text{aceite}$}{
        $s_i^{\text{acc}} \gets$ tentativa aceita; registrar confiança e justificativa
      }{
        $o \gets \Act{Executar}(\text{ação}, \text{argumentos})$
        \If{$\text{ação} = \text{avaliação de estratégia}$}{
          $sc^{\text{melhor}} \gets \min\{sc : (s, sc) \in \mathcal{H}_i\}$
          $\mathcal{H}_i \gets \mathcal{H}_i \cup \{(s_o, \text{score}_i(s_o))\}$
          \eIf{$\text{score}_i(s_o) < sc^{\text{melhor}}$}{
            $q \gets 0$
          }{
            $q \gets q + 1$
          }
        }
      }
    }
    $t \gets t + 1$
  }
  $s_i^{\star} \gets \arg\min\{sc : (s, sc) \in \mathcal{H}_i\}$
  \If{$s_i^{\text{acc}} \neq \text{nulo}$ \textbf{e} $s_i^{\text{acc}} \neq s_i^{\star}$}{
    registrar divergência
  }
  $f_i \gets$ regra de combinação de $s_i^{\star}$
  $\hat{\mathbf{y}}_i^{(\mathrm{comb})} \gets f_i(\hat{\mathbf{Y}}_i^{(\mathrm{teste})})$
  $\text{MAPE}_i, \text{SMAPE}_i, \text{RMSE}_i, \text{POCID}_i \gets \text{métricas}(\hat{\mathbf{y}}_i^{(\mathrm{comb})}, \mathbf{y}_i^{(\mathrm{teste})})$
}
\KwRet{$\{\hat{\mathbf{y}}_i^{(\mathrm{comb})}\}_{i=1}^{N}$; métricas; registros}
\end{algorithm}

\begin{algorithm}[H]
\caption{Mecanismo de pontuação de estratégias}
\label{alg:pontuar}
\Entrada{Série $i$; estratégia $s$}
\Saida{$\text{score}_i(s)$}

\For{$w \gets 1$ \KwTo $3$}{
  \If{$s$ seleciona o subconjunto a partir dos dados}{
    $\mathcal{M}_{i,s}^{(w)} \gets$ seleção sobre as janelas $\{1, 2, 3\} \setminus \{w\}$
  }
  \If{$s$ ajusta pesos a partir dos dados}{
    $\mathbf{p}_{i,s}^{(w)} \gets$ ajuste sobre as janelas $w' < w$
  }
  $\hat{\mathbf{y}}_{i,s}^{(w)} \gets \Act{Combinar}(s, \mathcal{M}_{i,s}^{(w)}, \mathbf{p}_{i,s}^{(w)})$
}
$\text{RMSE}_i(s), \text{SMAPE}_i(s), \text{MAPE}_i(s) \gets$ calculadas sobre as janelas concatenadas
$\text{POCID}_i(s) \gets \frac{1}{3}\sum_{w=1}^{3} \text{POCID}_i^{(w)}(s)$
$s_0 \gets$ média simples sobre $\mathcal{M}$
$\text{score}_i(s) \gets$ Equação~\eqref{eq:score}
\KwRet{$\text{score}_i(s)$}
\end{algorithm}

\begin{algorithm}[H]
\caption{Diagnóstico interpretativo}
\label{alg:diagnostico}
\Entrada{Perfil $\mathbf{z}_i$; avaliação do conjunto de preditores $T_i$}
\Saida{Diagnóstico $\mathcal{C}_i$}

\eIf{força de tendência alta \textbf{e} força sazonal alta}{
  $\text{regime} \gets \text{misto}$
}{
  \eIf{força sazonal alta}{
    $\text{regime} \gets \text{dominado por sazonalidade}$
  }{
    \eIf{força de tendência alta}{
      $\text{regime} \gets \text{dominado por tendência}$
    }{
      $\text{regime} \gets \text{ruidoso}$
    }
  }
}
\eIf{entropia espectral alta}{
  $\text{previsibilidade} \gets \text{baixa}$
}{
  \eIf{entropia espectral baixa}{
    $\text{previsibilidade} \gets \text{alta}$
  }{
    $\text{previsibilidade} \gets \text{média}$
  }
}
\eIf{$\tau$ baixo}{
  $\text{indicação} \gets \text{combinação robusta}$
}{
  \eIf{dispersão de erro alta \textbf{e} $\tau$ razoavelmente estável}{
    $\text{indicação} \gets \text{combinação seletiva}$
  }{
    \eIf{dispersão de erro alta}{
      $\text{indicação} \gets \text{combinação ponderada}$
    }{
      $\text{indicação} \gets \text{conjunto completo}$
    }
  }
}
$\text{riscos} \gets \emptyset$
\If{$\tau$ baixo}{$\text{riscos} \gets \text{riscos} \cup \{\text{instabilidade do ranking}\}$}
\If{há grupo de erro correlacionado}{$\text{riscos} \gets \text{riscos} \cup \{\text{redundância}\}$}
\If{previsibilidade baixa}{$\text{riscos} \gets \text{riscos} \cup \{\text{sobreajuste}\}$}
$\text{riscos} \gets \text{riscos} \cup \{\text{apenas três janelas}\}$
$\text{narrativa} \gets$ síntese textual das categorias
$\mathcal{C}_i \gets (\text{regime}, \text{previsibilidade}, \text{indicação}, \text{riscos}, \text{narrativa})$
\If{modelo de linguagem configurado}{
  $\mathcal{C}' \gets$ leitura do modelo restrita às quatro categorias
  \If{$\mathcal{C}'$ válido}{$\mathcal{C}_i \gets \mathcal{C}'$}
}
\KwRet{$\mathcal{C}_i$}
\end{algorithm}
\section{Preparação em nível de conjunto de dados} \label{sec:preparacao}

Antes que o agente atue sobre qualquer série individual, duas formas de conhecimento são construídas a partir do conjunto de dados $\mathcal{D}$ como um todo, sem que nenhum modelo de linguagem esteja envolvido: um meta-modelo que aprende a relação entre características de uma série e o erro de cada preditor $M_j$, e um levantamento de quais estratégias de combinação simples tendem a funcionar melhor naquele conjunto de dados especificamente. Ambas dependem de um mesmo mecanismo de pontuação, usado tanto aqui quanto, mais adiante, para avaliar as propostas do agente durante seu ciclo de decisão, descrito a seguir antes das duas formas de conhecimento propriamente ditas.

\subsection{Mecanismo de pontuação de estratégias} \label{subsec:pontuacao}

%Uma estratégia de combinação $s$ é definida por um subconjunto de preditores $\mathcal{M}_s \subseteq \mathcal{M}$ e por um método de agregação sobre as previsões desse subconjunto, podendo ser a média, a mediana, a média aparada, o DBA, ou uma média ponderada por um vetor de pesos $\mathbf{p}_s$. Aplicar $s$ à série $X_i$ na janela $w$ produz uma previsão combinada $\hat{\mathbf{y}}_{i,s}^{(w)} \in \mathbb{R}^h$.
Uma estratégia de combinação $s$ especifica uma regra de
seleção de preditores, um método de agregação e, quando
aplicável, uma regra de cálculo de pesos. O método de
agregação pode ser a média, mediana,
DBA ou média ponderada.

Para a série $i$ e a janela de validação $w$, essas regras
determinam um subconjunto não vazio
$\mathcal{M}_{i,s}^{(w)} \subseteq \mathcal{M}$ e, quando
aplicável, um vetor de pesos $\mathbf{p}_{i,s}^{(w)}$.
A aplicação da estratégia produz o vetor de previsões
combinadas
\[
    \hat{\mathbf{y}}_{i,s}^{(w)}
    =
    \bigl(\hat{y}_{i,s,k}^{(w)}\bigr)_{k=1}^{h}
    \in \mathbb{R}^{h}.
\]

Para avaliar $s$ na série $i$, os valores reais das três
janelas de validação são concatenados em um único vetor:
\[
    \mathbf{y}_i^{(1:3)}
    =
    \begin{bmatrix}
        \mathbf{y}_i^{(1)} \\
        \mathbf{y}_i^{(2)} \\
        \mathbf{y}_i^{(3)}
    \end{bmatrix}
    \in \mathbb{R}^{3h}.
\]
Analogamente, define-se
\[
    \hat{\mathbf{y}}_{i,s}^{(1:3)}
    =
    \begin{bmatrix}
        \hat{\mathbf{y}}_{i,s}^{(1)} \\
        \hat{\mathbf{y}}_{i,s}^{(2)} \\
        \hat{\mathbf{y}}_{i,s}^{(3)}
    \end{bmatrix}
    \in \mathbb{R}^{3h}.
\]

%Para avaliar $s$ sobre a série $X_i$, os valores reais das três janelas de validação são concatenados em um único vetor, denotado $\mathbf{y}_i^{(1:3)} \in \mathbb{R}^{3h}$, e o mesmo é feito com as previsões combinadas correspondentes, $\hat{\mathbf{y}}_{i,s}^{(1:3)} \in \mathbb{R}^{3h}$. 

As métricas RMSE, SMAPE e MAPE são calculadas sobre esse vetor concatenado, tratando as três janelas como uma sequência contínua de observações. O POCID constitui exceção a essa regra: por medir acerto de direção, uma grandeza que só é definida dentro de uma sequência temporal contínua, ele é calculado separadamente em cada uma das três janelas, e o resultado reportado é a média simples dos três valores.

Como as métricas resultantes têm escalas distintas entre si, elas não podem ser somadas diretamente sem que a de maior magnitude domine a soma. Para resolver isso, cada métrica de $s$ é expressa como uma razão contra a mesma métrica obtida pela estratégia de referência $s_0$, definida como a média simples sobre o conjunto completo $\mathcal{M}$, calculada com o mesmo procedimento. A pontuação de $s$ para a série $i$, denotada $\text{score}_i(s)$, é então:

\begin{equation}
\footnotesize
\text{score}_i(s) = \lambda_{\text{RMSE}} \frac{\text{RMSE}_i(s)}{\text{RMSE}_i(s_0)} + \lambda_{\text{SMAPE}} \frac{\text{SMAPE}_i(s)}{\text{SMAPE}_i(s_0)} + \lambda_{\text{MAPE}} \frac{\text{MAPE}_i(s)}{\text{MAPE}_i(s_0)} - \lambda_{\text{POCID}} \frac{\text{POCID}_i(s)}{100}
\label{eq:score}
\end{equation}
em que $(\lambda_{\text{RMSE}}, \lambda_{\text{SMAPE}}, \lambda_{\text{MAPE}}, \lambda_{\text{POCID}}) = (0{,}3,\ 0{,}3,\ 0{,}2,\ 0{,}2)$ neste trabalho. Como RMSE, SMAPE e MAPE são medidas de erro, em que valores menores são melhores, elas entram somando na Equação~\eqref{eq:score}; como o POCID é uma taxa de acerto, em que valores maiores são melhores, ele entra subtraindo, o que faz com que a convenção "quanto menor, melhor" valha para o \text{score} como um todo. Por construção, $s_0$ sempre obtém razão igual a $1$ em cada um dos três primeiros termos, de modo que $\text{score}_i(s_0)$ é uma constante determinada unicamente pelo seu próprio POCID; qualquer estratégia com $\text{score}_i(s) < \text{score}_i(s_0)$ superou a média simples naquela série.

Duas salvaguardas evitam que essa pontuação seja artificialmente otimista. A primeira diz respeito a estratégias cujos pesos $\mathbf{p}_{i,s}^{(w)}$ são ajustados a partir dos dados: o ajuste segue um esquema de janela expansiva, de modo que, ao pontuar a janela $w$, os pesos utilizados só podem ter sido calculados a partir das janelas anteriores a $w$, nunca a partir da própria janela sendo pontuada. A segunda diz respeito a estratégias cujo subconjunto $\mathcal{M}_s$ é selecionado a partir dos dados, como as baseadas nos preditores mais estáveis: essa seleção segue um esquema \textit{leave-one-out} entre janelas, de modo que o subconjunto usado para pontuar a janela $w$ é determinado a partir apenas das demais janelas, excluindo $w$. Essas duas salvaguardas não têm efeito sobre as estratégias clássicas, média, mediana e DBA aplicadas ao conjunto completo $\mathcal{M}$, já que essas não possuem peso ajustável nem subconjunto selecionado, não havendo, portanto, nenhum canal por onde a pontuação dessas três estratégias específicas pudesse vazar informação das janelas futuras.

\subsection{Meta-modelo entre séries}

O meta-modelo tem como objetivo aprender a relação entre as características de uma série e o desempenho de cada preditor de $\mathcal{M}$ sobre ela, de forma semelhante à realizada por \cite{fforma2019} no \textit{framework} FFORMA. Para cada histórico de valores da série $i$, i.e., $X_i \in \mathcal{D}$, é calculado um vetor de 26 características, composto por quatro medidas interpretáveis, a força de tendência, a força de sazonalidade, a entropia espectral e a autocorrelação de primeira ordem, e por 22 características adicionais obtidas do conjunto \textit{catch22} \citep{lubba2019catch22}. É importante notar que essas 26 características descrevem exclusivamente o comportamento da série; nenhuma delas deriva do erro histórico de qualquer preditor de $\mathcal{M}$.

O meta-modelo é treinado com o esquema \textit{leave-one-series-out}: para cada $X_i$, um modelo é ajustado utilizando as demais $N-1$ séries do conjunto, sem qualquer acesso às características ou aos erros obtidos a partir de $X_i$, e esse modelo é então o único empregado para gerar os pesos daquela série específica. Diferentemente de treinar um regressor separado para prever o erro de cada um dos preditores de $\mathcal{M}$ de forma independente, o meta-modelo é um único modelo que recebe as 26 características de uma série como entrada e emite $M$ valores brutos, transformados em pesos por meio de uma função \textit{softmax}. O treinamento otimiza diretamente o erro da combinação ponderada resultante desses pesos, e não o erro de previsão de cada peso individualmente, seguindo a mesma formulação do FFORMA.

\subsection{Prior de estratégias do conjunto de dados}

Além do meta-modelo, um segundo tipo de conhecimento entre séries é construído a partir de um conjunto fixo de estratégias de combinação simples, testadas em todas as séries de $\mathcal{D}$. Esse conjunto é composto por três estratégias clássicas aplicadas ao conjunto completo $\mathcal{M}$, a média, a mediana e o DBA, e por seis estratégias baseadas em estabilidade, a média e a média aparada calculadas sobre os $\kappa$ preditores de $\mathcal{M}$ com desempenho mais consistente entre as três janelas de validação, para $\kappa \in \{5, 7, 9\}$. Quando o meta-modelo descrito na subseção anterior está disponível para o conjunto de dados em questão, seus pesos constituem uma décima estratégia. A consistência empregada para selecionar os $\kappa$ preditores mais estáveis não é o menor erro médio entre as janelas, mas sim a estabilidade do \textit{ranking} de cada preditor ao longo delas: um preditor que ocupa, por exemplo, a segunda, a terceira e a quarta posição nas três janelas é considerado mais estável do que um que ocupa a primeira, a décima quinta e a terceira posição, ainda que os dois apresentem erro médio semelhante. Essa escolha evita utilizar a mesma métrica de validação tanto para selecionar o subconjunto de preditores quanto para, em seguida, pontuar esse mesmo subconjunto, o que inflacionaria artificialmente seu desempenho aparente.

Cada uma dessas estratégias fixas é avaliada em cada série de um conjunto de dados segundo o mecanismo de pontuação descrito na Seção~\ref{subsec:pontuacao}, obtendo $\text{score}_i(s)$. O prior de uma estratégia $s$ para o histórico da série $X_i$ é então definido como a média aritmética simples dessa pontuação nas demais séries do conjunto de dados:

\begin{equation}
\text{prior}_i(s) = \frac{1}{N-1} \sum_{i' \neq i} \text{score}_{i'}(s)
\label{eq:prior}
\end{equation}

O prior de cada estratégia fixa é disponibilizado ao agente como um indicativo do que costuma funcionar naquele conjunto de dados, sem que isso constitua uma regra a ser seguida. É importante destacar que essa média entre séries só é válida porque $\text{score}_{i'}(s)$, conforme definido na Equação~\eqref{eq:score}, já é uma razão adimensional em relação à própria estratégia de referência $s_0$ daquela série. Caso a agregação entre séries fosse feita diretamente sobre uma métrica em escala absoluta, como o RMSE, uma série de maior volume de vendas dominaria a média sem que isso indicasse, de fato, qual estratégia é mais adequada; ao operar sobre razões já normalizadas pela própria escala de cada série, o prior definido na Equação~\eqref{eq:prior} permanece comparável entre séries de magnitudes arbitrariamente distintas.




\section{Caracterização e diagnóstico por série} \label{sec:diagnostico}

Antes de o agente propor qualquer estratégia de combinação para uma série $i$, três formas de conhecimento específicas dessa série são calculadas: um perfil descritivo da própria série, uma avaliação do desempenho de cada preditor de $\mathcal{M}$ nas janelas de validação, e um diagnóstico interpretativo que traduz esses números em uma leitura qualitativa. As três, em conjunto, compõem o contexto inicial fornecido ao agente, antes de seu ciclo de decisão, descrito na Seção~\ref{sec:agente}.

\subsection{Perfil da série}

O perfil de $X_i$ é composto pelas mesmas 26 características empregadas no meta-modelo descrito na Seção~\ref{sec:preparacao}. Diferentemente do meta-modelo, cujo uso dessas características é interno ao cálculo dos pesos, aqui elas compõem parte do contexto textual fornecido diretamente ao agente, servindo de insumo tanto para sua própria leitura da série quanto para o diagnóstico interpretativo descrito adiante.

\subsection{Avaliação do conjunto de preditores}

Diferentemente do escore de estratégia definido na Seção~\ref{subsec:pontuacao}, que mede o desempenho de uma combinação inteira, esta etapa mede o desempenho de cada preditor $M_j \in \mathcal{M}$ isoladamente, usando duas grandezas distintas para finalidades diferentes.

A primeira é o RMSE de cada preditor $M_j$ em cada janela de validação $w$, calculado na escala original da série,

\begin{equation}
\text{RMSE}_{i,j}^{(w)} = \sqrt{\frac{1}{h} \sum_{k=1}^{h} \left(y_{i,k}^{(w)} - \hat{y}_{i,j,k}^{(w)}\right)^2},
\end{equation}
em que $y_{i,k}^{(w)}$ e $\hat{y}_{i,j,k}^{(w)}$ são, respectivamente, a componente $k$ do vetor de valores reais $\mathbf{y}_i^{(w)}$ e do vetor de previsões $\hat{\mathbf{y}}_{i,j}^{(w)}$, definidos na Seção~\ref{sec:formulacao}. A partir dele, são construídos dois relatórios: uma tabela que ranqueia os preditores por janela, mostrando os de melhor desempenho e um resumo dos demais; e uma medida de o quanto essa ordem se mantém estável entre as três janelas.

Essa estabilidade é medida pelo coeficiente de correlação de postos de Kendall, usualmente denotado $\tau$. De forma simples, esse coeficiente compara duas ordenações, neste caso o \textit{ranking} de preditores em uma janela de validação contra o \textit{ranking} na janela seguinte, e mede o quanto elas concordam entre si: para cada par de preditores, verifica-se se a ordem entre eles se manteve a mesma nas duas janelas, ou se inverteu. O coeficiente varia de $-1$, quando a ordem se inverte por completo, a $1$, quando a ordem se mantém idêntica, passando por $0$, quando não há relação alguma entre as duas ordenações. Um valor baixo de $\tau$ indica que a posição relativa de um preditor muda bastante de uma janela para outra, ou seja, que o melhor modelo em uma janela não é garantia de que continuará sendo o melhor na janela seguinte.

A segunda grandeza é o resíduo de cada preditor,
\[
    r_{i,j,k}^{(w)} = \hat{y}_{i,j,k}^{(w)} - y_{i,k}^{(w)},
\]
a diferença entre o valor previsto e o valor real, preservando o sinal dessa diferença. A partir dos resíduos, calcula-se a correlação entre cada par de preditores, permitindo identificar grupos que erram de forma parecida e na mesma direção ao longo das janelas. Preditores com erro fortemente correlacionado tendem a ser redundantes entre si: combiná-los agrega menos diversidade à combinação do que combinar preditores cujos erros não seguem esse mesmo padrão.

\subsection{Diagnóstico interpretativo}

A partir do perfil da série e da avaliação do conjunto de preditores, é construído um diagnóstico interpretativo, que traduz essas grandezas em quatro categorias textuais e uma narrativa curta. Por padrão, essa tradução é feita por um conjunto de regras determinísticas, sem envolvimento de nenhum modelo de linguagem; quando um modelo de linguagem é configurado para essa etapa, ele recebe as mesmas quatro categorias como um vocabulário fechado a ser preenchido por leitura das evidências, e a versão determinística permanece como reserva, empregada sempre que a resposta do modelo estiver ausente ou fora desse vocabulário.

O regime da série é determinado pela força de tendência e pela força de sazonalidade do perfil descrito nesta seção: quando as duas são altas, o regime é misto; quando só a sazonalidade é alta, é dominado por sazonalidade; quando só a tendência é alta, é dominado por tendência; e quando nenhuma das duas é alta, é ruidoso.

A previsibilidade é determinada pela entropia espectral da série, uma medida de o quanto seu comportamento se aproxima de um sinal aleatório: entropia alta indica previsibilidade baixa, entropia baixa indica previsibilidade alta, e valores intermediários, assim como a ausência dessa característica, indicam previsibilidade média.

A indicação de combinação nasce do cruzamento entre a estabilidade do \textit{ranking}, o $\tau$ descrito na subseção anterior, e a dispersão relativa de erro entre os preditores do conjunto, uma medida simples de o quanto o pior preditor do conjunto erra a mais do que o melhor deles. Quando o \textit{ranking} é pouco estável, a indicação é de combinação robusta, menos sensível à escolha exata de pesos; quando a dispersão é alta e o \textit{ranking} é razoavelmente estável, a indicação é de combinação seletiva, restrita a um subconjunto; quando só a dispersão é alta, a indicação é de combinação ponderada; fora desses casos, a indicação é de uso do conjunto completo, sem seleção ou ponderação adicional.

Por fim, é construída uma lista de riscos, a partir de condições independentes entre si, podendo conter mais de um item ao mesmo tempo: instabilidade do \textit{ranking} entre janelas; existência de algum grupo de preditores com erro correlacionado; previsibilidade baixa, sob o risco de que pesos ajustados a uma série próxima do ruído resultem em sobreajuste; e a limitação, presente em toda série avaliada por este método, de haver apenas três janelas de validação disponíveis para fundamentar qualquer estimativa de peso.

Essas quatro categorias são então combinadas em uma narrativa textual curta, entregue ao agente como parte do contexto inicial da série.


\section{Agente combinador} \label{sec:agente}
Antes da primeira iteração do agente, o histórico de tentativas da série já não está vazio: as mesmas dez estratégias fixas descritas na Seção~\ref{sec:preparacao}, avaliadas nas outras séries do conjunto de dados para compor o prior, são também avaliadas na própria série $i$, segundo o mecanismo de pontuação da Seção~\ref{subsec:pontuacao}, e seus resultados são inseridos no histórico antes de o agente atuar. Essas tentativas semeadas cumprem duas funções: servem de ponto de partida concreto para o raciocínio do agente, e garantem, por construção, um piso de desempenho, já que a estratégia final aplicada nunca pode ser pior que a melhor dessas dez, discutido na Seção~\ref{sec:aplicacao}.

Com o contexto da série já preparado, conforme descrito na Seção~\ref{sec:diagnostico}, o agente entra em ação para decidir a estratégia de combinação a ser empregada. Esta seção descreve o catálogo de ações sobre o qual o agente opera, o ciclo pelo qual ele raciocina e decide, e os critérios que encerram esse ciclo.

\subsection{Catálogo de ações}

O catálogo é composto por 24 ações determinísticas, organizadas em cinco categorias que correspondem às etapas naturais da tarefa de combinação: diagnóstico, seleção de pool, cálculo de pesos, combinação e validação. Cada ação é implementada como uma função convencional, sem qualquer participação de modelo de linguagem em sua execução, e o agente interage com elas exclusivamente pelo nome da ação e pelos argumentos fornecidos em sua \textit{Action Input}, descrita na Seção~\ref{sec:agente}.

\textbf{Ações de diagnóstico.} Produzem as evidências sobre a série e sobre o conjunto de preditores $\mathcal{M}$, parte das quais já compõe o contexto inicial fornecido ao agente, descrito na Seção~\ref{sec:diagnostico}, mas que também podem ser recalculadas pelo agente ao longo do ciclo de decisão, com parâmetros distintos dos usados nesse contexto inicial. As ações disponíveis nessa categoria e suas respectivas descrições são apresentadas na Tabela~\ref{tab:acoes_descricao}.

\begin{table}[H]
\centering
\caption{Ações de diagnóstico e suas respectivas descrições.}
\label{tab:acoes_descricao}
\begin{tabular}{@{}p{5cm}p{9cm}@{}}
\toprule
\textbf{Ação} & \textbf{Descrição} \\
\midrule
Perfil da série & Retorna comprimento, sazonalidade, estacionariedade, presença de \textit{outliers} e características do conjunto \textit{catch22}. \\
Resumo da decomposição & Retorna a proporção da variância da série explicada por tendência, sazonalidade e resíduo. \\
Resumo de erro dos preditores & Retorna a tabela de RMSE por preditor, ranqueada, limitada aos de melhor desempenho e a um resumo dos demais. \\
Estabilidade do \textit{ranking} & Retorna o coeficiente de estabilidade do \textit{ranking} entre janelas e identifica quais preditores mais mudam de posição. \\
Correlação de erro & Retorna grupos de preditores cujos resíduos são fortemente correlacionados entre si, indicando redundância. \\
Teste de Diebold-Mariano & Compara dois preditores nomeados, com a correção para amostras pequenas de \citet{harvey1997testing}. \\
\bottomrule
\end{tabular}
\end{table}

\textbf{Ações de seleção de pool.} Produzem um subconjunto de $\mathcal{M}$, referenciado por um identificador que pode ser usado como argumento de ações posteriores, sem que o agente precise enumerar os preditores individualmente. A Tabela~\ref{tab:acoes_selecao_pool} apresenta as ações disponíveis para selecionar e refinar esses subconjuntos.

\begin{table}[H]
\centering
\caption{Ações de seleção de pool de preditores.}
\label{tab:acoes_selecao_pool}
\begin{tabular}{@{}p{5cm}p{9cm}@{}}
\toprule
\textbf{Ação} & \textbf{Descrição} \\
\midrule
Seleção pelos $\kappa$ melhores &
Seleciona os $\kappa$ preditores de menor erro nas janelas indicadas. \\
Seleção pelos mais estáveis &
Seleciona os $\kappa$ preditores mais consistentes entre janelas, combinando a posição média no \textit{ranking} com sua variação. \\
Poda de redundância &
Remove de um subconjunto os preditores redundantes identificados pela ação de correlação de erro, mantendo apenas o de menor erro em cada grupo. \\
\bottomrule
\end{tabular}
\end{table}

\textbf{Ações de cálculo de pesos.} Produzem um vetor de pesos sobre um subconjunto de preditores, também referenciado por um identificador, nunca pelos valores numéricos em si. Conforme apresentado na Tabela~\ref{tab:acoes_pesos}, as ações dessa categoria utilizam diferentes critérios para determinar a contribuição de cada preditor na combinação.

\begin{table}[H]
\centering
\caption{Ações de cálculo de pesos e suas respectivas descrições.}
\label{tab:acoes_pesos}
\begin{tabular}{@{}p{5cm}p{9cm}@{}}
\toprule
\textbf{Ação} & \textbf{Descrição} \\
\midrule
Peso pelo erro inverso & Pesos proporcionais ao inverso do erro de cada preditor, com um parâmetro que permite aproximar o resultado do peso uniforme. \\
Peso por \textit{softmax} do erro & Pesos proporcionais ao \textit{softmax} do erro negativo, a formulação empregada pelo ADE \citep{ade}. \\
Peso pela tendência do erro & Pesos que consideram para onde o erro de cada preditor está evoluindo ao longo das janelas, não apenas sua média; calculados, por padrão, a partir do erro absoluto médio (MAE) de cada preditor, para que um único ponto de erro elevado não domine a direção da tendência estimada. \\
Peso por mínimos quadrados & Pesos obtidos por mínimos quadrados, na formulação de \citet{granger1984}, projetados sobre o simplex de probabilidade por padrão. \\
Peso por características da série & Pesos de um meta-modelo ajustado especificamente para a série, a partir de suas características, no espírito do FFORMA \citep{fforma2019}; recorre à ação de \textit{softmax} do erro quando há poucas janelas disponíveis. \\
Peso pelo meta-modelo entre séries & Pesos do meta-modelo descrito na Seção~\ref{sec:preparacao}, treinado com as demais séries do conjunto de dados. \\
\bottomrule
\end{tabular}
\end{table}

\textbf{Ações de combinação.} Produzem a previsão combinada propriamente dita, a partir de um subconjunto de preditores e, quando aplicável, de um vetor de pesos. As diferentes estratégias de combinação disponibilizadas ao agente estão descritas na Tabela~\ref{tab:acoes_combinacao}.

\begin{table}[H]
\centering
\caption{Ações de combinação de previsões.}
\label{tab:acoes_combinacao}
\begin{tabular}{@{}p{5cm}p{9cm}@{}}
\toprule
\textbf{Ação} & \textbf{Descrição} \\
\midrule
Combinação por média & Média simples entre os preditores do subconjunto. \\
Combinação por mediana & Mediana entre os preditores do subconjunto, mais robusta a um preditor com previsão discrepante. \\
Combinação por média aparada & Descarta uma proporção das previsões mais extremas em cada extremidade antes de calcular a média. \\
Combinação ponderada & Média ponderada, a partir de um identificador de pesos já calculado por uma ação da categoria anterior. \\
Combinação por DBA & Alinha temporalmente as previsões, por meio de \textit{Dynamic Time Warping Barycenter Averaging}, antes de combiná-las. \\
Seleção de um único preditor & Adota um único preditor nomeado, sem qualquer combinação. \\
\bottomrule
\end{tabular}
\end{table}

\textbf{Ações de validação.} Não produzem uma nova previsão, mas avaliam e organizam o que já foi produzido pelas ações anteriores. A Tabela~\ref{tab:acoes_validacao} apresenta as ações utilizadas para avaliar estratégias, verificar a plausibilidade das previsões e consultar os resultados de tentativas anteriores.

\begin{table}[H]
\centering
\caption{Ações de validação e acompanhamento das estratégias.}
\label{tab:acoes_validacao}
\begin{tabular}{@{}p{5cm}p{9cm}@{}}
\toprule
\textbf{Ação} & \textbf{Descrição} \\
\midrule
Avaliação da estratégia & Monta a estratégia especificada, executa o procedimento de \textit{backtest} descrito na Seção~\ref{subsec:pontuacao} e insere o resultado no histórico de tentativas da série; é a única ação por meio da qual uma estratégia passa a integrar esse histórico. \\
Verificação de sanidade & Compara a previsão de teste de uma estratégia contra os limites históricos da série, emitindo um alerta sem impedir a continuidade do ciclo. \\
Listagem do histórico & Retorna o histórico de tentativas já registradas para a série, ranqueado do melhor para o pior resultado. \\
\bottomrule
\end{tabular}
\end{table}

Além dessas 24 ações, existe uma operação adicional, o aceite, que não integra o catálogo no mesmo sentido das demais: em vez de ser executada e devolver uma observação para uma nova iteração, ela encerra o ciclo de decisão da série, formalizando o critério de aceite descrito na Seção~\ref{sec:agente}. O aceite é acompanhado de uma confiança autodeclarada pelo agente e de uma justificativa para a estratégia escolhida, ambas registradas junto ao histórico da série para fins de análise posterior.
\subsection{Ciclo de raciocínio e ação}

O agente é implementado sobre um modelo de linguagem de peso aberto, executado localmente, operando segundo o paradigma ReAct discutido na Seção~\ref{subsec:react}. A cada iteração, o agente recebe um conjunto de informações textuais sobre o estado atual da decisão e deve produzir, em resposta, exatamente três linhas: um \textit{Thought}, em que verbaliza seu raciocínio sobre o que fazer a seguir; uma \textit{Action}, que identifica qual ação do catálogo deseja invocar; e uma \textit{Action Input}, que especifica os argumentos dessa ação. Nenhum conteúdo além dessas três linhas é aceito como parte da decisão do agente.

O texto fornecido ao agente a cada iteração combina uma parte fixa, que não se altera entre iterações da mesma série, e uma parte que se atualiza a cada novo turno. A parte fixa estabelece o papel do agente, o catálogo de ações disponível, o formato de saída exigido, e um conjunto de diretrizes práticas de raciocínio, entre as quais duas se destacam: a orientação de que pesos distribuídos de forma quase uniforme sobre o conjunto de preditores é um indício de que o problema não está na forma como os pesos são calculados, mas na composição do subconjunto de preditores escolhido, e a exigência de que qualquer estratégia aceita seja justificada em termos de características observáveis da série, e não apenas pelo valor numérico do erro obtido. A parte variável inclui o número de iterações restantes, o perfil e o diagnóstico da série descritos na Seção~\ref{sec:diagnostico}, o prior de estratégias do conjunto de dados descrito na Seção~\ref{sec:preparacao}, explicitamente apresentado como um indicativo e não como uma regra a seguir, e o histórico de tentativas já realizadas nessa série, ordenado da melhor para a pior, cada uma acompanhada da justificativa fornecida pelo agente no momento em que foi proposta.

A resposta do agente passa por uma etapa de validação antes de qualquer ação ser de fato executada. Caso o nome da ação solicitada não exista no catálogo, ou caso os argumentos fornecidos não sejam válidos para a ação solicitada, uma mensagem de erro estruturada é devolvida ao agente, que pode corrigir sua próxima resposta a partir dela. Duas categorias de falha na comunicação com o modelo, uma resposta vazia ou um erro de conexão, são tratadas separadamente: como não representam uma decisão do agente, são reexecutadas automaticamente, até um número limitado de tentativas, sem consumir uma iteração do orçamento de decisão. Já uma resposta que viola o formato de três linhas esperado é tratada como uma decisão mal formada do próprio agente, consumindo uma iteração e devolvendo um lembrete do formato exigido.

\subsection{Critérios de encerramento}

O ciclo de decisão de uma série é encerrado por uma de três condições. A primeira é o aceite explícito, quando o agente decide, entre as estratégias já testadas no histórico daquela série, qual delas deve ser adotada como final. A segunda é a estagnação, quando um número consecutivo de tentativas propostas pelo agente deixa de superar a melhor estratégia já registrada no histórico, levando o sistema a encerrar o ciclo mesmo sem um aceite explícito. A terceira é o esgotamento do orçamento de iterações, um limite máximo de tentativas por série, que atua como rede de segurança independentemente do comportamento do agente. Em qualquer um dos três casos, a estratégia levada adiante para a etapa seguinte, descrita na Seção~\ref{sec:aplicacao}, é sempre a de melhor pontuação entre todas as tentativas do histórico daquela série, segundo o mecanismo descrito na Seção~\ref{subsec:pontuacao}, incluindo as estratégias fixas testadas antes de o agente atuar.






\section{Aplicação da estratégia e rastreabilidade} \label{sec:aplicacao}

\subsection{Aplicação na janela de teste}

Ao término do ciclo de decisão descrito na Seção~\ref{sec:agente}, por qualquer um dos três critérios de encerramento, a estratégia levada adiante é sempre aquela de melhor pontuação entre todas as tentativas registradas no histórico da série, segundo o mecanismo descrito na Seção~\ref{subsec:pontuacao}. O histórico é ordenado exclusivamente por essa pontuação, do menor para o maior valor, sem qualquer margem mínima de diferença entre a primeira e a segunda colocada: uma tentativa com pontuação menor assume a liderança do histórico independentemente de quão pequena seja essa diferença. Isso significa que, quando o agente aceita explicitamente uma estratégia, essa aceitação só é levada adiante se a estratégia aceita for, de fato, a de melhor pontuação do histórico; caso o agente aceite uma estratégia que não ocupa essa posição, é a estratégia de melhor pontuação, e não a aceita pelo agente, que é aplicada, e essa divergência é registrada para fins de rastreabilidade, discutidos na Seção~\ref{subsec:auditoria}.

A estratégia vencedora é então aplicada, pela primeira e única vez, às previsões de teste dos preditores de $\mathcal{M}$: o mesmo subconjunto $\mathcal{M}_s$ e o mesmo método de agregação, sejam eles a média, a mediana, a média aparada, o DBA ou a combinação ponderada por $w_s$, conforme definidos na Seção~\ref{subsec:pontuacao}, são recalculados agora sobre as previsões da janela de teste em vez das janelas de validação, produzindo a previsão final $\hat{y}_i^{(\text{combinada})}$. O cálculo das métricas de desempenho sobre essa previsão é deliberadamente distinto do mecanismo de pontuação empregado durante o ciclo de decisão: não há normalização contra a estratégia de referência $s_0$, nem ponderação entre métricas. São calculadas as quatro métricas descritas na Seção~\ref{sec:metricas}, MAPE, SMAPE, RMSE e POCID, diretamente entre a previsão combinada e os valores reais da janela de teste, sem qualquer transformação adicional. Essas quatro métricas, e não a pontuação interna que orientou a decisão, são o resultado reportado deste método.

% \subsection{Procedência} \label{subsec:auditoria}

% Além da aplicação da estratégia vencedora, um mecanismo de procedência avalia e registra a qualidade da decisão tomada, permitindo tanto a análise crítica dos resultados quanto a reprodução integral do processo por trás de cada previsão: ele verifica que toda previsão, peso ou métrica que compõe o resultado final de uma série de fato se origina da execução real de uma ação do catálogo, e não de algum outro caminho. Essa garantia decorre, em primeiro lugar, de uma restrição estrutural da implementação: uma tentativa só pode ser registrada por uma única rotina interna, por sua vez acionada exclusivamente por duas vias, a semeadura das estratégias fixas descrita na Seção~\ref{sec:preparacao} ou a ação de avaliação de estratégia descrita na Seção~\ref{sec:agente}. Como a saída do agente é sempre texto, o nome de uma ação e seus argumentos, e nunca um valor numérico diretamente, essa saída só pode disparar uma dessas duas vias, cada uma das quais recalcula toda a evidência a partir das previsões originais dos preditores de $\mathcal{M}$; não existe, portanto, um caminho pelo qual um valor pudesse ser inserido no resultado final sem ter sido de fato computado dessa forma.

% Para tornar essa garantia auditável, e não apenas presumida pelo desenho da implementação, dois registros são mantidos e verificados ao final do processamento de cada série. O primeiro é um registro de todas as chamadas de ação realizadas, bem-sucedidas ou não, a partir do qual se verifica que o número de chamadas de avaliação de estratégia bem-sucedidas é sempre maior ou igual ao número de tentativas de origem do agente presentes no histórico daquela série, sem exigir igualdade estrita, já que reenviar uma estratégia já testada anteriormente gera uma nova chamada registrada, mas não uma nova tentativa. O segundo é uma verificação de forma sobre a evidência associada a cada tentativa do histórico: o número de resíduos registrados deve corresponder ao produto entre o número de janelas de validação e o horizonte de previsão, e o número de pontuações por janela deve corresponder ao número de janelas, confirmando que essa evidência tem a forma esperada de um resultado de fato calculado sobre as três janelas, ainda que não confirme diretamente cada valor individual. Essa verificação é executada automaticamente ao final de cada série, e seu resultado é registrado permanentemente para fins de análise; caso ela indique alguma inconsistência, contudo, a série não deixa de ser processada, seguindo o mesmo princípio de degradação sem interrupção adotado em outras partes do método, como o retorno ao diagnóstico determinístico quando um modelo de linguagem configurado para essa etapa falha, discutido na Seção~\ref{sec:diagnostico}.





\section{Considerações finais} \label{sec:consideracoes-finais-proposta}

Este capítulo apresentou o CREST, método desenvolvido para a combinação dinâmica de previsões de séries temporais, que delega a um agente baseado em modelo de linguagem a decisão de qual estratégia de combinação empregar para cada série, em substituição a uma regra fixa aplicada indistintamente a todas elas.

O método parte de duas restrições identificadas na formulação do problema: a ausência de uma estratégia de combinação universalmente superior entre séries e entre conjuntos de dados, e a escassez de evidência disponível, apenas três janelas de validação, para fundamentar diretamente a estimação de pesos ótimos. Em resposta a essas restrições, a decisão de combinação foi estruturada em duas macro-fases: uma preparação em nível de conjunto de dados, que constrói conhecimento entre séries por meio de um meta-modelo e de um prior de estratégias, sem envolvimento de nenhum modelo de linguagem; e uma decisão por série, na qual um agente itera sobre um catálogo fechado de ações determinísticas, validando cada proposta por meio de um mecanismo de pontuação comum antes de aceitá-la.

A decisão de projeto central do método, a de que o agente nunca produz diretamente um valor numérico, orientou o desenho de cada componente descrito neste capítulo: o catálogo de ações restringe o espaço de decisão do agente a operações verificáveis; a estratégia final aplicada corresponde sempre à melhor pontuação de todo o histórico, e não necessariamente à escolha do agente; e o mecanismo de procedência torna essa decisão passível de verificação posterior, em vez de dependente unicamente do julgamento do próprio agente.

O capítulo seguinte apresenta os resultados obtidos pela aplicação deste método aos conjuntos de dados descritos, comparando seu desempenho ao dos modelos individuais, das estratégias de combinação estática e dos métodos de combinação por meta-aprendizagem revisados no Capítulo~\ref{chp:revisao}.